\documentclass[main.tex]{subfiles}


\begin{document}

%from pagenumber to up
% \phantomsection 
% \hypertarget{toc}{}

 

 

 
   

\section{Сила Лоренца}


\textbf{Сила Лоренца} ($F_\Lambda$ [Н])~--- это сила, действующая на \emph{движущийся} заряд со стороны магнитного поля:
\begin{equation}
    F_\Lambda=qBv\sin\alpha,
\end{equation}
где $q$~--- величина заряда, $B$~--- индукция магнитного поля, $v$~--- скорость заряда, $\alpha$~--- угол между направлениями скорости~$v$ и поля~$B$.

Пусть заряд влетает в пространство между полюсами магнита (рис.\,\ref{fig:p155}).


    \begin{figure}[H]
    \centering
\includegraphics[]{el/39bFL_pic1.pdf}%
\caption{Отклонение движения заряда в магнитном поле}
\label{fig:p155}
\end{figure}


Положительный заряд (красный шар), движущийся со скоростью~$\vec v$, при попадании в поле магнита (оно сосредоточено практически между его полюсами) отклоняется от первоначальнго направления движения под действием силы Лоренца (для сравнения серой штриховой линией показана траектория движения незаряженной частицы). Вот правило для определения направления этой силы.
%
\begin{law}
    \textbf{Правило левой руки.} Если расположить левую руку так, чтобы вытянутые четыре пальца указывали направление скорости \emph{положительного} заряда, а линии поля входили в ладонь, то отведенный перпендикулярно большой палец (лежащий в одной плоскости с остальными пальцами) укажет направление силы Лоренца (если заряд отрицательный, то четыре пальца устанавливают \emph{против} его скорости).
\end{law}
%
% Рисунок~\,\ref{fig:p156} иллюстрирует применение этого правила.  
%
    \begin{figure}[H]
    \centering
\includegraphics[]{el/39bFL_pic2.pdf}%
\caption{Правило левой руки}
\label{fig:p156}
\end{figure}






























 
\end{document}